\scanpage{88}

% top, coproduct
Coproduit \inserted{d'un graphe} avec un doublet $\bullet\quad\bullet$ donne, pour un graphe discret, les \unsure{bigraphes|graphes} bidiscrets.

% upper, doublets
Coproduit de deux doublets $(\underset a\bullet\ \underset b\bullet)\ (\underset{a'}\bullet\ \underset{b'}\bullet)$

% upper, crossed pairs
\diagram{\begin{tikzcd}[ampersand replacement=\&,row sep=small,column sep=small]
(ba')\arrow[dr,dash] \& (bb')\arrow[dl,dash]\\
(aa') \& (ab')
\end{tikzcd}}

% upper middle, segments
\unsure{Coproduit|produit} de deux segments $(\underset{a}{\bullet}\!-\!\underset{b}{\bullet})\quad(\underset{a'}{\bullet}\!-\!\underset{b'}{\bullet})$

% upper middle, square 1
\diagram{\textcircled{1}\quad\begin{tikzcd}[ampersand replacement=\&,row sep=small,column sep=small]
(ba')\arrow[r,dash]\arrow[d,dash]\arrow[dr,dash,draw=gray] \& (bb')\arrow[d,dash]\arrow[dl,dash,draw=gray]\\
(aa')\arrow[r,dash] \& (ab')
\arrow[from=1-1,to=2-2,phantom,"\deleted{\times}"]
\end{tikzcd}\qquad carré}

% middle, opposite and bracketed passage
\deleted{\unsure{dual}} opposé du précédent. [Le \uline{produit} de \deleted{\unsure{un doublet} $\underset{a}{\bullet}\ \underset b\bullet$ et d'un segment $\underset{a'}\bullet\!-\!\underset{b'}\bullet$} \inserted{deux doublets est l'ens à quatre él. discret} \deleted{\unsure{alors par}} le coproduit de deux segments est donc le graphe chaotique à $4$ él. -- ces deux \unsure{résultats|doublets|semble} \unsure{peu|purs} \unsure{enthousiasmants|contrairement|intéressants}\ldots]

% middle below fold
Produit d'un doublet \inserted{$(a\bullet\ \bullet b)$} et d'un segment \inserted{$a'\bullet\!-\!\bullet b'$}

% lower middle, two horizontal segments
\diagram{\begin{tikzcd}[ampersand replacement=\&,row sep=small,column sep=small]
(ba')\arrow[r,dash] \& (bb')\\ (aa')\arrow[r,dash] \& (ab')
\end{tikzcd},}

% lower, opposite
son opposé est le coproduit du segment $(\overset a\bullet\!-\!\overset b\bullet)$ et du doublet $(a'\bullet\ \bullet b')$

% lower, square 2
\diagram{\textcircled{2}\quad\begin{tikzcd}[ampersand replacement=\&,row sep=small,column sep=small]
(ba')\arrow[d,dash]\arrow[dr,dash] \& (bb')\arrow[d,dash]\arrow[dl,dash,crossing over]\\
(aa') \& (ab')
\end{tikzcd}\qquad\uline{carré}}

% foot, continuation
Donc sur le produit de deux ens de cardinal $2$, on trouve $3$ structures de carré naturelles distinctes, \deleted{\unsure{comme}} savoir la structure \uline{produit} de deux \uline{segments}

% archive sheet number: 87
